\documentclass[12pt,a4paper]{article}

% Encodage et langue
\usepackage[utf8]{inputenc}
\usepackage[T1]{fontenc}
\usepackage[french]{babel}

% Mise en page et utilitaires
\usepackage{geometry}
\usepackage{graphicx}
\usepackage{float}
\usepackage{array}
\usepackage{setspace}
\usepackage{hyperref}
\usepackage{xcolor}
\usepackage{caption}

\geometry{
    top=1.8cm,
    bottom=1.8cm,
    left=2cm,
    right=2cm
}

\setstretch{1.15}
\setlength{\parskip}{0.35em}
\setlength{\parindent}{0pt}

\hypersetup{
    colorlinks=true,
    linkcolor=black,
    urlcolor=blue
}

\begin{document}

% ---------------- Page de garde ----------------
\begin{titlepage}
    \thispagestyle{empty}

    \noindent
    \begin{minipage}[t]{0.48\textwidth}
        \IfFileExists{bdiologo.png}
            {\includegraphics[height=2cm,width=1.6cm]{bdiologo.png}}
            {\fbox{\rule{0pt}{2cm}\hspace{1cm}BDIO\hspace{1cm}}}
    \end{minipage}
    \hfill
    \begin{minipage}[t]{0.48\textwidth}
        \raggedleft
        \IfFileExists{ensamlogo.png}
            {\includegraphics[height=1.6cm]{ensamlogo.png}}
            {\fbox{\rule{0pt}{1.6cm}\hspace{1cm}ENSAM\hspace{1cm}}}
    \end{minipage}

    \vspace{3.2cm}

    \begin{center}
        {\small Projet JEE \& BigData}

        \vspace{0.35cm}

        {\LARGE \bfseries SkillBridge}

        \vspace{0.45cm}

        {\Large Description du sujet}

        \vspace{0.8cm}

        {\normalsize
        SkillBridge est une plateforme éducative orientée projet qui part d'une idée concrète
        proposée par l'utilisateur pour l'aider à comprendre quoi apprendre, dans quel ordre,
        et vers quels cours s'orienter. L'objectif est de transformer un besoin souvent flou en
        un parcours d'apprentissage clair, pertinent et directement lié au projet à réaliser.}
    \end{center}

    \vspace{4.8cm}

    \noindent
    \begin{minipage}[t]{0.42\textwidth}
        {\small Réalisé par :}

        \vspace{0.35cm}

        {\small Omar Elkhali}

        \vspace{0.25cm}

        {\small Ayoub Moutik}

        \vspace{0.25cm}

        {\small Jamila Ouzzane}
    \end{minipage}
    \hfill
    \begin{minipage}[t]{0.32\textwidth}
        {\small Encadré par :}

        \vspace{0.35cm}

        {\small Pr. EL Faquih Loubna}

        \vspace{0.25cm}

        {\small Pr. HIRCHOUA Bader}

        \vspace{0.25cm}

        {\small Pr. SEKKATE Sara}
    \end{minipage}

    \vfill

    \begin{center}
        {\small AU : 25-26}
    \end{center}

\end{titlepage}

% ---------------- Sommaire ----------------
\renewcommand{\contentsname}{Sommaire}
\tableofcontents
\newpage

% ---------------- Rapport ----------------
\section{Introduction}
\textbf{SkillBridge} est une plateforme éducative qui part d'une idée de projet proposée par l'apprenant pour la transformer en parcours d'apprentissage. Au lieu de demander simplement ``quel cours veux-tu suivre ?'', la plateforme pose une question plus utile : \textbf{``qu'est-ce que tu veux construire ?''}. À partir de cette intention, elle cherche à faire le lien entre le projet visé, les compétences à acquérir et les cours les plus pertinents.

Le projet s'inscrit dans le domaine de \textbf{l'éducation numérique}. Il répond à une difficulté simple : beaucoup d'étudiants savent ce qu'ils veulent réaliser, mais ne savent pas clairement quoi apprendre pour y arriver. Face à des ressources nombreuses mais dispersées, l'objectif de SkillBridge est de transformer un besoin concret en un parcours d'apprentissage cohérent. La problématique du projet peut donc se résumer ainsi : \textbf{comment partir d'une idée de projet et orienter l'utilisateur vers les compétences et les cours les plus pertinents ?}

\section{Architecture globale du système}
L'organisation générale de \textbf{SkillBridge} repose sur deux ensembles complémentaires : une application principale destinée aux interactions avec l'utilisateur, et un pipeline Big Data chargé de collecter, stocker et analyser les données produites par la plateforme.

Deux acteurs interviennent dans ce système : \textbf{l'utilisateur} et \textbf{l'administrateur}. L'utilisateur peut créer un compte, soumettre une idée de projet, consulter les recommandations proposées, enregistrer des cours et suivre sa progression. L'administrateur, de son côté, gère le catalogue de cours, les catégories, les fournisseurs ainsi que les données générales nécessaires au bon fonctionnement de la plateforme.

Sur le plan fonctionnel, l'application prend en charge l'authentification, la gestion des utilisateurs, la gestion du catalogue, la saisie d'idées de projet, la recommandation de cours, l'enregistrement de favoris et le suivi de progression. En parallèle, le pipeline Big Data a pour rôle de centraliser les événements d'usage et de produire des données d'analyse qui pourront enrichir la compréhension du comportement utilisateur et améliorer la pertinence des recommandations.

Une évolution envisagée consiste à intégrer un wrapper d'API d'intelligence artificielle, comme Claude ou ChatGPT, dans la partie frontend. Son rôle serait d'aider à interpréter la demande formulée par l'utilisateur et à construire une première arborescence de compétences à partir du prompt saisi. La même logique pourrait aussi être appliquée aux cours collectés afin de produire des informations structurées avant leur traitement dans le pipeline Big Data.

Cette séparation permet de distinguer clairement la partie opérationnelle, centrée sur l'expérience utilisateur, de la partie analytique, centrée sur l'exploitation des données.

\begin{figure}[H]
    \centering
    \includegraphics[width=0.95\textwidth]{Image/Capture d'écran 2026-04-09 220804.png}
    \caption{Architecture globale de SkillBridge}
\end{figure}

Le schéma d'architecture met ainsi en évidence les acteurs du système, la séparation entre l'application principale et le pipeline Big Data, ainsi que la circulation des données entre les différentes parties de la solution.

\section{Sources de données envisagées}
Pour alimenter le catalogue de cours, plusieurs sources ont déjà été identifiées. À ce stade, les principales ressources retenues sont les jeux de données suivants :

\begin{itemize}
    \item Kaggle --- \textit{EdX, Coursera, and Udemy Course Data} ;
    \item Kaggle --- \textit{Online Course Dataset (EdX, Udacity, Coursera)} ;
    \item Kaggle --- \textit{Global Online Tech Courses Dataset}.
\end{itemize}

Ces jeux de données offrent une première base pour récupérer des métadonnées de cours, comparer plusieurs plateformes et enrichir le catalogue de départ. En complément, il est aussi prévu de développer un outil capable de scraper automatiquement des cours gratuits afin d'ajouter de nouvelles ressources de manière progressive.

\section{Modélisation du système}

L'entité centrale est \textbf{User}, qui représente l'utilisateur de la plateforme. Chaque utilisateur possède un \textbf{Role}, ce qui permet de distinguer les droits d'un utilisateur simple de ceux d'un administrateur. Un utilisateur peut soumettre plusieurs \textbf{ProjectIdea}, enregistrer plusieurs \textbf{SavedCourse} et suivre son avancement dans différents cours à travers l'entité \textbf{Progress}.

L'entité \textbf{Course} représente un cours accessible depuis la plateforme. Chaque cours est rattaché à une \textbf{Category} pour le classement thématique, et à un \textbf{Provider}, qui représente l'organisme ou la source du contenu. Les notions ou compétences couvertes par un cours sont représentées par l'entité \textbf{Skill}.

Lorsqu'un utilisateur soumet une idée de projet, celle-ci peut être liée à plusieurs compétences. À partir de ce lien, le système peut produire une ou plusieurs \textbf{Recommendation}, orientées vers les cours les plus pertinents. Cette structure reste simple, claire et suffisamment extensible pour les prochaines phases du projet.

\begin{figure}[H]
    \centering
    \includegraphics[width=0.86\textwidth]{Image/image.png}
    \caption{Diagramme de classes UML de SkillBridge}
\end{figure}

\section{Choix technologiques}
\begin{itemize}
    \item \textbf{React} : pour une interface dynamique et fluide.
    \item \textbf{Spring Boot} : pour structurer le backend rapidement et proprement.
    \item \textbf{Spring Security} : pour gérer l'authentification et les droits d'accès.
    \item \textbf{PostgreSQL} : pour stocker les données relationnelles du système.
    \item \textbf{Sqoop} : pour échanger les données entre la base relationnelle et l'environnement Big Data.
    \item \textbf{Flume} : pour collecter les événements et les logs.
    \item \textbf{HDFS} : pour le stockage distribué des données massives.
    \item \textbf{HBase} : pour conserver certaines données analytiques exploitables rapidement.
    \item \textbf{MapReduce} et \textbf{Hive} : pour les traitements batch et les requêtes analytiques.
    \item \textbf{Python} : pour les traitements complémentaires sur les données et le texte.
\end{itemize}

\end{document}
